% Lo que garantiza cada axioma sobre los conjuntos de indiferencia
% (fusión de los gráficos 2.2a-f de Segura en seis paneles, en el orden de los axiomas del tema)
% Autor: Víctor Gutiérrez Marcos
\documentclass[tikz,border=4pt]{standalone}
% ==========================================================================
% tikz-tcee.tex — Estilos comunes de los gráficos TikZ del temario
% Autor: Víctor Gutiérrez Marcos
%
% Cada gráfico es un fichero standalone en fuentes/<tema>/graficos/:
%
%   \documentclass[tikz,border=4pt]{standalone}
%   % ==========================================================================
% tikz-tcee.tex — Estilos comunes de los gráficos TikZ del temario
% Autor: Víctor Gutiérrez Marcos
%
% Cada gráfico es un fichero standalone en fuentes/<tema>/graficos/:
%
%   \documentclass[tikz,border=4pt]{standalone}
%   % ==========================================================================
% tikz-tcee.tex — Estilos comunes de los gráficos TikZ del temario
% Autor: Víctor Gutiérrez Marcos
%
% Cada gráfico es un fichero standalone en fuentes/<tema>/graficos/:
%
%   \documentclass[tikz,border=4pt]{standalone}
%   \input{../../../latex/tikz-tcee}
%   \begin{document}
%   \begin{tikzpicture}[tcee]
%     \ejes{Q}{P}{6}{5}
%     \draw[curva1] (0.5,4.5) -- (5,0.8) node[etiqueta,right] {$D$};
%     \capa{2}{\draw[curva2] ...;}   % en el vídeo aparece en el paso 2
%   \end{tikzpicture}
%   \end{document}
%
% Paleta: la de la web (morado, dorado, salvia) más tres colores de apoyo
% distinguibles en blanco y negro por el trazo.
% ==========================================================================
\usepackage[T1]{fontenc}
\usepackage{newpxtext,newpxmath}
\usepackage{xcolor}
\usepackage{tikz,pgfplots}
\pgfplotsset{compat=1.18}
\usetikzlibrary{arrows.meta,calc,positioning,patterns,decorations.pathreplacing,intersections,shapes.misc,backgrounds}

\definecolor{tceemorado}{HTML}{5F2987}
\definecolor{tceedorado}{HTML}{B8860B}
\definecolor{tceeverde}{HTML}{3C7D3A}
\definecolor{tceeazul}{HTML}{1F5F99}
\definecolor{tceerojo}{HTML}{B22222}
\definecolor{tceegris}{HTML}{767171}
\definecolor{tceesalvia}{HTML}{E2EFD9}

% Capas para el vídeo: \capa{n}{…} dibuja su contenido solo a partir del paso n.
% En el PDF, la web y Word se ve todo; generar-video.py compila el gráfico una
% vez por paso con \def\tceepaso{k} para que cada elemento aparezca al narrarlo.
\providecommand{\tceepaso}{1000}
\newcommand{\capa}[2]{\ifnum#1>\tceepaso\relax\else#2\fi}

\tikzset{
  tcee/.style={line cap=round,line join=round,font=\small,>={Stealth[length=2.2mm]}},
  eje/.style={->,thick,black},
  curva1/.style={very thick,tceemorado},
  curva2/.style={very thick,tceeazul},
  curva3/.style={very thick,tceeverde},
  curva4/.style={very thick,tceedorado},
  curvanueva/.style={very thick,tceerojo},
  desplazada/.style={very thick,dashed},
  guia/.style={thin,dashed,tceegris},
  punto/.style={circle,fill=black,inner sep=1.3pt},
  etiqueta/.style={font=\small,inner sep=2pt},
  flecha/.style={->,thick,tceerojo},
  area/.style={fill=tceemorado,fill opacity=0.18,draw=none},
  area2/.style={fill=tceedorado,fill opacity=0.22,draw=none},
}

% \ejes{etiqueta x}{etiqueta y}{largo x}{alto y}
\newcommand{\ejes}[4]{%
  \draw[eje] (0,0) -- (#3,0) node[below left=2pt and 0pt] {#1};
  \draw[eje] (0,0) -- (0,#4) node[left] {#2};
  \node[below left] at (0,0) {$0$};}
% \proyeccion{(x,y)}{etq x}{etq y}: líneas guía a los ejes con etiquetas
\newcommand{\proyeccion}[3]{%
  \draw[guia] let \p1=#1 in (\x1,0) node[below] {#2} |- (0,\y1) node[left] {#3};}

%   \begin{document}
%   \begin{tikzpicture}[tcee]
%     \ejes{Q}{P}{6}{5}
%     \draw[curva1] (0.5,4.5) -- (5,0.8) node[etiqueta,right] {$D$};
%     \capa{2}{\draw[curva2] ...;}   % en el vídeo aparece en el paso 2
%   \end{tikzpicture}
%   \end{document}
%
% Paleta: la de la web (morado, dorado, salvia) más tres colores de apoyo
% distinguibles en blanco y negro por el trazo.
% ==========================================================================
\usepackage[T1]{fontenc}
\usepackage{newpxtext,newpxmath}
\usepackage{xcolor}
\usepackage{tikz,pgfplots}
\pgfplotsset{compat=1.18}
\usetikzlibrary{arrows.meta,calc,positioning,patterns,decorations.pathreplacing,intersections,shapes.misc,backgrounds}

\definecolor{tceemorado}{HTML}{5F2987}
\definecolor{tceedorado}{HTML}{B8860B}
\definecolor{tceeverde}{HTML}{3C7D3A}
\definecolor{tceeazul}{HTML}{1F5F99}
\definecolor{tceerojo}{HTML}{B22222}
\definecolor{tceegris}{HTML}{767171}
\definecolor{tceesalvia}{HTML}{E2EFD9}

% Capas para el vídeo: \capa{n}{…} dibuja su contenido solo a partir del paso n.
% En el PDF, la web y Word se ve todo; generar-video.py compila el gráfico una
% vez por paso con \def\tceepaso{k} para que cada elemento aparezca al narrarlo.
\providecommand{\tceepaso}{1000}
\newcommand{\capa}[2]{\ifnum#1>\tceepaso\relax\else#2\fi}

\tikzset{
  tcee/.style={line cap=round,line join=round,font=\small,>={Stealth[length=2.2mm]}},
  eje/.style={->,thick,black},
  curva1/.style={very thick,tceemorado},
  curva2/.style={very thick,tceeazul},
  curva3/.style={very thick,tceeverde},
  curva4/.style={very thick,tceedorado},
  curvanueva/.style={very thick,tceerojo},
  desplazada/.style={very thick,dashed},
  guia/.style={thin,dashed,tceegris},
  punto/.style={circle,fill=black,inner sep=1.3pt},
  etiqueta/.style={font=\small,inner sep=2pt},
  flecha/.style={->,thick,tceerojo},
  area/.style={fill=tceemorado,fill opacity=0.18,draw=none},
  area2/.style={fill=tceedorado,fill opacity=0.22,draw=none},
}

% \ejes{etiqueta x}{etiqueta y}{largo x}{alto y}
\newcommand{\ejes}[4]{%
  \draw[eje] (0,0) -- (#3,0) node[below left=2pt and 0pt] {#1};
  \draw[eje] (0,0) -- (0,#4) node[left] {#2};
  \node[below left] at (0,0) {$0$};}
% \proyeccion{(x,y)}{etq x}{etq y}: líneas guía a los ejes con etiquetas
\newcommand{\proyeccion}[3]{%
  \draw[guia] let \p1=#1 in (\x1,0) node[below] {#2} |- (0,\y1) node[left] {#3};}

%   \begin{document}
%   \begin{tikzpicture}[tcee]
%     \ejes{Q}{P}{6}{5}
%     \draw[curva1] (0.5,4.5) -- (5,0.8) node[etiqueta,right] {$D$};
%     \capa{2}{\draw[curva2] ...;}   % en el vídeo aparece en el paso 2
%   \end{tikzpicture}
%   \end{document}
%
% Paleta: la de la web (morado, dorado, salvia) más tres colores de apoyo
% distinguibles en blanco y negro por el trazo.
% ==========================================================================
\usepackage[T1]{fontenc}
\usepackage{newpxtext,newpxmath}
\usepackage{xcolor}
\usepackage{tikz,pgfplots}
\pgfplotsset{compat=1.18}
\usetikzlibrary{arrows.meta,calc,positioning,patterns,decorations.pathreplacing,intersections,shapes.misc,backgrounds}

\definecolor{tceemorado}{HTML}{5F2987}
\definecolor{tceedorado}{HTML}{B8860B}
\definecolor{tceeverde}{HTML}{3C7D3A}
\definecolor{tceeazul}{HTML}{1F5F99}
\definecolor{tceerojo}{HTML}{B22222}
\definecolor{tceegris}{HTML}{767171}
\definecolor{tceesalvia}{HTML}{E2EFD9}

% Capas para el vídeo: \capa{n}{…} dibuja su contenido solo a partir del paso n.
% En el PDF, la web y Word se ve todo; generar-video.py compila el gráfico una
% vez por paso con \def\tceepaso{k} para que cada elemento aparezca al narrarlo.
\providecommand{\tceepaso}{1000}
\newcommand{\capa}[2]{\ifnum#1>\tceepaso\relax\else#2\fi}

\tikzset{
  tcee/.style={line cap=round,line join=round,font=\small,>={Stealth[length=2.2mm]}},
  eje/.style={->,thick,black},
  curva1/.style={very thick,tceemorado},
  curva2/.style={very thick,tceeazul},
  curva3/.style={very thick,tceeverde},
  curva4/.style={very thick,tceedorado},
  curvanueva/.style={very thick,tceerojo},
  desplazada/.style={very thick,dashed},
  guia/.style={thin,dashed,tceegris},
  punto/.style={circle,fill=black,inner sep=1.3pt},
  etiqueta/.style={font=\small,inner sep=2pt},
  flecha/.style={->,thick,tceerojo},
  area/.style={fill=tceemorado,fill opacity=0.18,draw=none},
  area2/.style={fill=tceedorado,fill opacity=0.22,draw=none},
}

% \ejes{etiqueta x}{etiqueta y}{largo x}{alto y}
\newcommand{\ejes}[4]{%
  \draw[eje] (0,0) -- (#3,0) node[below left=2pt and 0pt] {#1};
  \draw[eje] (0,0) -- (0,#4) node[left] {#2};
  \node[below left] at (0,0) {$0$};}
% \proyeccion{(x,y)}{etq x}{etq y}: líneas guía a los ejes con etiquetas
\newcommand{\proyeccion}[3]{%
  \draw[guia] let \p1=#1 in (\x1,0) node[below] {#2} |- (0,\y1) node[left] {#3};}

\usetikzlibrary{shapes.geometric}
\begin{document}
\begin{tikzpicture}[tcee,
    titulo/.style={font=\small\bfseries,align=center,anchor=base},
    pref/.style={draw=tceemorado,thick,regular polygon,regular polygon sides=3,inner sep=1.4pt}]
  \newcommand{\ejesp}{\draw[eje] (0,0) -- (4.3,0) node[below left=2pt and 0pt] {$x_1$};
    \draw[eje] (0,0) -- (0,3.8) node[left] {$x_2$};
    \node[below left] at (0,0) {$0$};}
  % ---------- 1: preorden completo (axiomas 1-3) ----------
  \capa{1}{
  \begin{scope}[shift={(0,0)}]
    \node[titulo] at (2.1,4.1) {Axiomas 1-3: preorden completo};
    \ejesp
    \foreach \p in {(0.5,3.3),(0.75,2.85),(0.45,2.45),(1.2,2.25),(2.3,2.6),(3.0,3.3),(0.45,1.35),(1.6,1.35),(1.0,0.85),(2.65,1.0),(2.1,0.55)}
      \node[punto] at \p {};
    \foreach \p in {(1.75,3.45),(1.25,2.95),(3.15,2.75),(0.8,1.95),(1.25,1.65),(2.05,1.65),(2.95,1.35),(3.35,1.35),(0.65,0.75),(1.75,0.6),(2.95,0.65),(3.6,2.05)}
      \node[pref] at \p {};
    \node[punto,fill=tceerojo] (x0) at (2.4,1.95) {};
    \node[etiqueta,above right,text=tceerojo] at (x0) {$x^0$};
    \node[font=\footnotesize,align=left,anchor=north] at (2.1,-0.45)
      {\tikz\node[punto]{}; $\sim x^0$ \quad \tikz\node[pref]{}; $\succ x^0$ \quad resto: $\prec x^0$};
  \end{scope}}
  % ---------- 2: continuidad (axioma 4) ----------
  \capa{2}{
  \begin{scope}[shift={(5.6,0)}]
    \node[titulo] at (2.1,4.1) {+ 4. Continuidad};
    \ejesp
    \fill[area] plot[smooth] coordinates {(0,3.5) (0.8,2.8) (1.5,2.95) (2.3,2.2) (3.2,1.7) (4.0,1.6)}
      -- (4.0,0.45) -- plot[smooth] coordinates {(3.2,0.6) (2.4,0.75) (1.6,1.5) (0.9,1.45) (0,2.4)} -- cycle;
    \draw[curva1] plot[smooth] coordinates {(0,3.5) (0.8,2.8) (1.5,2.95) (2.3,2.2) (3.2,1.7) (4.0,1.6)};
    \draw[curva1] plot[smooth] coordinates {(0,2.4) (0.9,1.45) (1.6,1.5) (2.4,0.75) (3.2,0.6) (4.0,0.45)};
    \node[punto] (x0) at (1.9,1.95) {};
    \node[etiqueta,right] at (x0) {$x^0$};
    \node[etiqueta] at (3.1,1.15) {$I(x^0)$};
    \node[etiqueta] at (2.9,3.1) {$B(x^0)$};
    \node[etiqueta] at (0.75,0.6) {$H(x^0)$};
    \node[font=\footnotesize,align=center,anchor=north] at (2.1,-0.45) {conjuntos de contorno cerrados;\\ la indiferencia puede ser «gruesa»};
  \end{scope}}
  % ---------- 3: monotonía estricta (axioma 5) ----------
  \capa{3}{
  \begin{scope}[shift={(11.2,0)}]
    \node[titulo] at (2.1,4.1) {+ 5. Monotonía estricta};
    \ejesp
    \draw[curva1] plot[domain=0:3.8,samples=80,smooth] (\x,{3.4-0.85*\x+0.25*sin(2.5*\x r)})
      node[etiqueta,above right] {$I(x^0)$};
    \node[punto] (x0) at (0.8,{3.4-0.85*0.8+0.25*sin(2 r)}) {};
    \node[etiqueta,above right] at (x0) {$x^0$};
    \node[etiqueta] at (3.3,3.4) {$B(x^0)$};
    \node[etiqueta] at (0.8,0.8) {$H(x^0)$};
    \draw[flecha] (1.8,1.7) -- ++(0.7,0.7) node[etiqueta,right,text=tceerojo] {preferidas};
    \node[font=\footnotesize,align=center,anchor=north] at (2.1,-0.45) {curva «fina» y estrictamente\\ decreciente};
  \end{scope}}
  % ---------- 4: convexidad (axioma 6) ----------
  \capa{4}{
  \begin{scope}[shift={(0,-5.8)}]
    \node[titulo] at (2.1,4.1) {+ 6. Convexidad};
    \ejesp
    \draw[curva1] (0,3.3) -- (1.2,1.5) -- (3.9,0.25) node[etiqueta,above right] {$I(x^0)$};
    \node[punto] (x0) at (0.6,2.4) {};
    \node[etiqueta,right] at (x0) {$x^0$};
    \node[punto] (x1) at (1.2,1.5) {};
    \node[etiqueta,above right] at (x1) {$x^1$};
    \node[etiqueta] at (3.0,2.6) {$B(x^0)$};
    \node[etiqueta] at (0.75,0.6) {$H(x^0)$};
    \node[font=\footnotesize,align=center,anchor=north] at (2.1,-0.45) {convexa, pero con tramos\\ rectos y vértices};
  \end{scope}}
  % ---------- 5: convexidad estricta (axioma 6') ----------
  \capa{5}{
  \begin{scope}[shift={(5.6,-5.8)}]
    \node[titulo] at (2.1,4.1) {+ 6$'$. Convexidad estricta};
    \ejesp
    \draw[curva1] plot[domain=0.3:1.3,samples=40,smooth] (\x,{1.2+0.8*(1.3-\x)^2+1.4*(1.3-\x)})
      -- plot[domain=1.3:3.8,samples=60,smooth] (\x,{0.42+0.1154*(3.9-\x)^2})
      node[etiqueta,above] {$I(x^0)$};
    \node[punto] (x0) at (0.7,{1.2+0.8*0.36+1.4*0.6}) {};
    \node[etiqueta,right] at (x0) {$x^0$};
    \node[punto] (x1) at (1.3,1.2) {};
    \node[etiqueta,above right] at (x1) {$x^1$};
    \node[etiqueta] at (3.0,2.6) {$B(x^0)$};
    \node[etiqueta] at (0.75,0.6) {$H(x^0)$};
    \node[font=\footnotesize,align=center,anchor=north] at (2.1,-0.45) {sin tramos rectos;\\ aún puede haber vértices};
  \end{scope}}
  % ---------- 6: diferenciabilidad (axioma 7) ----------
  \capa{6}{
  \begin{scope}[shift={(11.2,-5.8)}]
    \node[titulo] at (2.1,4.1) {+ 7. Diferenciabilidad};
    \ejesp
    \draw[curva1] plot[domain=0.35:3.8,samples=80,smooth] (\x,{1.2/\x}) node[etiqueta,above] {$I(x^0)$};
    % tangente en x0 = (0.7, 1.2/0.7): pendiente -1.2/0.49
    \draw[guia,solid,tceegris,thick] plot[domain=0.3:1.1] (\x,{1.7143-2.449*(\x-0.7)});
    \node[punto] (x0) at (0.7,1.7143) {};
    \node[etiqueta,right] at (x0) {$x^0$};
    \node[etiqueta] at (3.0,2.6) {$B(x^0)$};
    \node[etiqueta] at (0.55,0.4) {$H(x^0)$};
    \node[font=\footnotesize,align=center,anchor=north] at (2.1,-0.45) {curva suave: pendiente (RMS)\\ única en cada punto};
  \end{scope}}
\end{tikzpicture}
\end{document}
